% !TeX root = ../main.tex
\section{How well it works}
\label{sec:results}

\subsection{The benchmark}

\code{eval_tone} (an ignored test in \file{tone/eval.rs}) uses EGDB test items:
each has the DI of a performance, renders of the same DI through three
amplifier simulations (Marshall, Mesa, Plexi: equipment \emph{outside} the app's
chain) and labelled notes. Per item there are four kinds of target:
\begin{description}[style=nextline,leftmargin=1.2cm]
\item[chain] The DI through a random template chain followed by a random
  cabinet-like EQ (a high-pass near 80\Hz, a roll-off above 3--7\kHz, six
  smooth bumps of up to $\pm4\dB$). The true settings are known, so parameter
  error can be measured. Random settings: a random voicing, each parameter
  uniform in 0.1--0.9, and the compressor, drive and reverb each switched off
  with probability 0.5, 0.4 and 0.5. In the working tree half of these targets
  also get a delay of 150--700\ms{} \wip.
\item[Marshall, Mesa, Plexi] The EGDB renders: can the app's chain imitate gear
  it does not contain?
\end{description}

Each target is estimated by five methods, all ending in the same kind of chain:
\begin{enumerate}
\item \textbf{dry DI}: no processing at all (a floor to beat).
\item \textbf{best factory preset}: the closest of the 20 presets, as is.
\item \textbf{preset + match}: that preset plus the match IR.
\item \textbf{search + match, from notes}: the full estimator, with the stand-in
  dry signal from the plucked-string model playing the \emph{labelled} notes.
  This is the real use case, except that the notes are perfect.
\item \textbf{search + match, true DI}: the full estimator given the real DI.
  An upper bound: what the search achieves when the dry signal is right.
\end{enumerate}

\begin{keyidea}[: how each estimate is scored]
Whatever dry signal the estimator used, the \emph{true} DI is rendered through
its estimated chain and compared with the target by the plain log-band distance
(no EQ allowance, Section~\ref{sec:app-distance}). This asks the question the
user cares about: if I play this performance through the proposed preset, how
far is it from the record? It also makes the five methods comparable.
\end{keyidea}

\subsection{Results}

Committed run (\code{7db5015}): 12 EGDB items, 12\,s excerpts, about 15\,s per
estimate on 28 threads \measured.

\begin{table}[h]
\centering
\caption{Mean log-band distance (dB, lower is better) between the true DI
rendered through each estimate and the target, 12 EGDB items \measured.}
\label{tab:tone-results}
\footnotesize
\setlength{\tabcolsep}{4pt}
\begin{tabular}{@{}lccccc@{}}
\toprule
Target & Dry DI & Best preset & Preset + match & Search, notes & Search, true DI \\
\midrule
Random app chain + cab EQ & 6.60 & 9.29 & 6.64 & \textbf{5.30} & 0.66 \\
EGDB Marshall & 7.52 & 7.55 & 6.84 & \textbf{6.49} & 3.11 \\
EGDB Mesa & 7.94 & 8.90 & 6.52 & \textbf{5.83} & 2.67 \\
EGDB Plexi & 7.61 & 8.16 & 7.63 & \textbf{6.26} & 2.32 \\
\bottomrule
\end{tabular}
\end{table}

\paragraph{What the table says.}
\begin{itemize}
\item \textbf{The estimator works, modestly.} Search from notes beats every
  baseline on average in all four groups, and in 7--8 of the 12 items per group.
  The gain over the best baseline is 0.35--1.35\dB.
\item \textbf{The match filter matters.} A factory preset alone is often
  \emph{worse} than no processing (9.29 against 6.60 on chain targets), because
  its colour is wrong; with the match IR it gets close to or better than the dry
  DI.
\item \textbf{The search itself is good; the stand-in dry signal is the
  bottleneck.} With the true DI the search gets within 0.66\dB{} of targets made
  with the app's own chain, essentially solving them, and to 2.3--3.1\dB{} of
  real amp simulations. From the notes it stays at 5.3--6.5\dB. The difference
  (3.2--4.6\dB) is the cost of the plucked-string stand-in, which is still
  9.9\dB{} from real DI even with an EQ allowed.
\item \textbf{Settings are not recovered, and that is expected.} On chain targets
  the mean parameter error is 0.41 on the 0--1 scale and the amp voicing is right
  in 7 of 12 items. Several settings sound alike (Section~\ref{sec:comparison}),
  so the search finds \emph{a} chain that sounds close, not \emph{the} chain.
\end{itemize}

\begin{worked}[: what a decibel of distance means]
The score is an average absolute error per band and frame. 6\dB{} means each
band of each frame is, on average, off by about a factor of two in amplitude
(20\,log$_{10}$\,2 $\approx$ 6\dB). It is a coarse measure: it cannot tell a
small error everywhere from a large error in a few bands, and it has not been
checked against listeners for these sounds.
\end{worked}

\begin{caution}[: what this benchmark does not yet test]
\begin{itemize}
\item \textbf{Perfect notes}: the stand-in plays labelled notes. With
  transcribed notes, missed and extra notes will add error \expected.
\item \textbf{Isolated guitar}: no target is a separated stem. The working tree
  adds separated-stem targets (\code{GVA_TONE_STEMS}), scored against the clean
  render they came from \wip.
\item \textbf{Small sample}: 12 items. Differences of a few tenths of a dB
  between methods are within item-to-item variation.
\item \textbf{The score is not validated}: no listening test yet. The benchmark
  writes WAVs (\code{GVA_TONE_OUT}) for listening.
\item \textbf{Only simulated amps} for the gear-outside-the-chain case; the
  working tree adds Guitar-TECHS miked real amplifiers \wip.
\end{itemize}
\end{caution}

\subsection{Being measured in the working tree}

At the time of writing, the benchmark in the working tree \wip{} adds: the
sampler stand-in as a sixth method (\code{GVA_TONE_SAMPLES}); delay-time accuracy
(right within 5\% when the target has a delay, ``no delay'' when it has none);
separated-stem and Guitar-TECHS miked-amp targets; a higher search floor
(\code{GVA_TONE_FLOOR}) so the search ignores quiet bands where a stand-in
differs most from a real DI; and \code{GVA_TONE_PLUCK=0} to skip the pluck
stand-in. No numbers from it are recorded yet; when they are, this section and
Table~\ref{tab:tone-results} should be updated.
